\section{Writing It Down Somewhere Else}
\label{sec:ch17-writing-it-down-somewhere-else}

\index{ownership!where the record can actually live|(}%
Section~\ref{sec:ch17-handing-a-field-over} ends on a directive I would not
ship. There is one I would, and it is the only thing in the whole set that
lets a subgraph change what a client sees of somebody else's field.

\index{inaccessible@\code{"@inaccessible}!one document takes a field off the menu}%
Mark the shared \code{abstract} \code{@inaccessible} in the Ratings document
alone. That composes at exit zero, and the config that comes out has grown a
key it did not have before.

\index{router execution config!a third key, and a second pair}%
Chapter~\ref{ch:composition} already pulled two schemas out of this file, the
composed \code{graphqlSchema} and the \code{serviceSdl} that keeps each
subgraph's own document whole beside it. This is a third key and a different
pair. \code{graphqlSchema} is what the router plans against and what
chapter~\ref{ch:breaking-changes} diffed, and until this edit it was the only
client-facing document any config in this book carried. Now
\code{graphqlClientSchema} sits beside it, nine lines shorter. Two of the nine
are the hidden field and its description. The rest is a client schema being a
client schema: the directive definitions go, and so do the descriptions inside
the \code{ApplyPolicy} enum that chapter~\ref{ch:entities-authorization} put
in the supergraph without anybody asking for it. WunderGraph's own contracts
page names the pair and says which is which~\autocite{cosmocontracts}: the
router schema keeps such fields because query planning still needs them, and
the client schema is what introspection and clients are given.

\index{inaccessible@\code{"@inaccessible}!hiding is not a routing question}%
So one team can remove a field from the public schema by editing its own
document. That is the exception to
section~\ref{sec:ch17-what-the-config-calls-a-subgraph}, which said every
directive here answers a routing question, and it is an exception for a
reason: hiding a field is not a routing question, the composer has no route to
lose over it, and so it does that one without asking anybody either.

\subsection{The Governance Products, and What They Need}

\index{tag@\code{"@tag}!what contracts are built out of}%
\code{@tag} is the other half of that machinery and it is inert on its own. It
labels schema elements, and a \emph{contract} is a named filter over those
labels: the elements tagged for removal come out of the composed schema, with
\code{@inaccessible} applied to them on the way, and what a given audience is
served is the smaller document that results~\autocite{cosmocontracts}. Cosmo
creates one with \code{wgc contract create}. Apollo builds the same feature out
of the same two directives, and describes the division between them in one
sentence: \code{@tag} on a type or field says whether to include or exclude it
from a contract schema, and \code{@inaccessible} is what you reach for instead
when you want it gone from every schema the graph serves, contracts and source
alike~\autocite{apollocontracts}.

\index{control plane!contracts and proposals both need one}%
Both are control-plane features, in the sense
chapter~\ref{ch:breaking-changes} established for the schema check: the
subcommands authenticate through \code{COSMO\_API\_KEY}, which for Cosmo Cloud
arrives with the account~\autocite{cosmocli}. So does \code{wgc proposal},
which is Cosmo's other governance product in the same CLI: a proposal is a
schema change published for review before it is published for real, and the
review runs five checks against it, composition, breaking changes, lint,
pruning, and the operation impact that needs the recorded client traffic
chapter~\ref{ch:breaking-changes} priced~\autocite{cosmoproposals}. That is
the shape of the thing this
chapter has been describing the absence of, and it is sold rather than
shipped, for the same reason chapter~\ref{ch:breaking-changes} could not run a
check: the record it works from lives on a server.

\index{ownership!access control is not ownership}%
The nearest thing either vendor has to an owner is access control. Cosmo
scopes roles at four levels, organization, namespace, federated graph and
subgraph, and at the subgraph level the roles are admin, publisher, checker
and viewer~\autocite{cosmogrouprules}. That is genuinely useful and it is a
different statement. Publisher is a permission, not a claim about whose type
it is, and a team holding it on the Ratings subgraph has, by that grant alone,
everything it needs to make both of the edits in
section~\ref{sec:ch17-the-field-nobody-approved}.

\subsection{What This Costs Teams Who Are Not Me}

\index{ownership!the collision two teams find at the gateway}%
The failure has a published shape. Isha Talegaonkar, an engineer at Walmart
Global Tech, wrote up the schema governance tooling her team built for exactly
this and named the problem in one line~\autocite{talegaonkar2022governance}:

\begin{quote}
Teams may end up creating the same \emph{Type} as another team, which can lead
to conflicts after integration in the federated gateway.
\end{quote}

\index{ownership!nobody finds out until the gateway}%
And the cost, a few lines further down the same post, is when they find out:
\enquote{developers go through the entire Software Development Life Cycle
(SDLC) before noticing the collisions}. Her team's answer was a linter with a
collision checker in it, run in the development cycle rather than at the
gateway, reporting to what the post calls the schema owner. That is a term
their tool has and the toolchain under it does not.

\index{ownership!the owner who cannot fix their own errors}%
The other shape is not a collision at all, and Netflix published it while
describing something else. Alex Hutter, Falguni Jhaveri and Senthil Sayeebaba
wrote about making a federated graph searchable, over a graph whose \enquote{three
core entities} are each \enquote{owned by separate engineering teams}, and the
paragraph worth reading twice is in their list of open
problems~\autocite{hutter2022search}:

\begin{quote}
The owner of the index is often required to follow up with other domain teams
regarding errors in related entities and be in the unenviable position of not
being able to do much to resolve the issues independently.
\end{quote}

\index{ownership!responsibility without the ability to act}%
That is ownership as it actually lands. Somebody owns the thing the user
complains about, and the data behind half of it belongs to teams they can
only write to. It is the position chapter~\ref{ch:blame-routing} is about
arriving one layer up, and no directive in this chapter would help them.

\subsection{What I Would Do}

\index{ownership!my judgment on what it costs}%
Ownership is what kills federation projects, ahead of every mechanical failure
in part III of this book. That is my judgment and it cites nobody: I looked
for a published account of a team arguing about who owns a field in a
federated graph, across eleven engineering blogs and the GraphQLConf
schedules, and the two above are the closest anybody comes. Nobody writes that
one up. The reasoning is what this chapter measured. Every mechanical failure
here has a message, a status code or a count attached to it, and somebody can
be handed it. An ownership failure produces exit zero and a config, and the
first artifact of it is a meeting.

\index{ownership!the document is a file, so use what files have}%
So write the owner where the document is. Every edit in this chapter went
through one file, in one repository, and a review on that file is the only
gate in the entire chain that a human reads. A code-ownership file over
\code{schema.graphql} does more for this than any directive, because it is the
one mechanism that fires on the thing that actually happened: somebody changed
a document.

\index{ownership!a field on their type is a change to their type}%
Treat a field you add to somebody else's type as a change to their type,
because it is one and the composer will not tell them.
Chapter~\ref{ch:second-third-subgraph} gave Ratings three fields on
\code{Session} and that was the right design; what makes it right is that the
fields are ratings data, hanging off a key. A field on their type whose data
is theirs is a duplicate with a directive on it, which is
section~\ref{sec:ch17-handing-a-field-over}'s \code{@provides} argument again.

\index{ownership!the shared root needs one owner by convention}%
Give the root type an owner by convention, since it will not have one any
other way. Chapter~\ref{ch:breaking-changes} counted seven root fields from
four services on one \code{Query} that none of them owns, in a graph with one
author and no clients. Whoever that owner is, their job is the naming and the
argument about whether a caller needs a new root field or an argument on an
existing one, which is the choice chapter~\ref{ch:breaking-changes} traced a
schema's whole shape back to.

\index{override@\code{"@override}!and delete it when it is done}%
And when a field does change hands, do it in the order
section~\ref{sec:ch17-handing-a-field-over} gives, and then go back and delete
the directive. It composes forever whether or not it still means anything, so
the only thing that will ever remove it is somebody remembering.

\index{ownership!the question the next chapter asks at run time}%
All of that assumes you can say whose field it is once somebody asks. That is
the easy version of the question, because a schema change has a document and a
commit behind it and somebody typed both. The hard version arrives at run time
and it is the next chapter's: the graph is slow, one request crossed four
services, and the composed schema says which of them answers each field and
nothing whatever about which of them took the time.
\index{ownership!where the record can actually live|)}%
